% ==============================================================================
% Macros for Relational Separation Logic and RTL Compiler Verification
% ==============================================================================

% Operational Semantics and Reductions
\newcommand{\step}[3]{#1 \vdash #2 \twoheadrightarrow #3}
\newcommand{\steps}[3]{#1 \vdash #2 \twoheadrightarrow^{\ast} #3}
\newcommand{\psteps}[3]{#1 \vdash #2 \twoheadrightarrow^{+} #3}
\newcommand{\nstep}[4]{#1 \vdash #2 \underset{#3}{\twoheadrightarrow} #4}

\newcommand{\isval}[1]{\textsf{is\_value}\left(#1\right)}
\newcommand{\bfinal}[3]{\textsf{both\_final}_{#1}\left(#2, #3\right)}
\newcommand{\canp}[2]{\textsf{can\_progress}_{#1}\left(#2\right)}
\newcommand{\stuck}[2]{\textsf{stuck}_{#1}\left(#2\right)}
\newcommand{\Some}[1]{\textsf{Some}\left(#1\right)}
\newcommand{\None}{\textsf{None}}

% Behavioral Predicates
\newcommand{\Spec}{\textsf{Spec}}
\newcommand{\Safe}[1]{\textsf{Safe}(#1)}
\newcommand{\Term}[2]{\textsf{Terminating}(#1, #2)}
\newcommand{\Div}{\textsf{Diverging}}
\newcommand{\Undef}{\textsf{Undefined}}

% Data Structures
\newcommand{\List}[1]{\textsf{List}\left(#1\right)}
\newcommand{\Stream}[1]{\textsf{Stream}\left(#1\right)}

% Simulation Relations
\newcommand{\isim}[3]{#1 \lesssim #2 \; \{#3\}}
\newcommand{\fsim}[5]{#1 \leftindex[I]_{#2}\lesssim_{#3} #4 \; \{#5\}}
\newcommand{\incirc}[1]{\raisebox{.5pt}{\textcircled{\raisebox{-.9pt}{#1}}}}

% Value Types and Constructors
\newcommand{\val}{\textsf{val}}
\newcommand{\loc}{\textsf{loc}}
\newcommand{\VInt}[1]{\textsf{VInt}\left(#1\right)}
\newcommand{\VBool}[1]{\textsf{VBool}\left(#1\right)}
\newcommand{\VPtr}[1]{\textsf{VPtr}\left(#1\right)}
\newcommand{\VUndef}{\textsf{VUndef}}

% Memory Cell Types and Constructors
\newcommand{\cell}{\textsf{cell}}
\newcommand{\Allocated}[1]{\textsf{Allocated}\left(#1\right)}
\newcommand{\Freed}{\textsf{Freed}}
\newcommand{\memory}{\textsf{memory}}

% Relational Value Similarity
\newcommand{\sameval}[3]{#2 \mathrel{\bowtie}_{#1} #3}

% Regmap Operations
\newcommand{\getreg}[2]{\ensuremath{#2[#1]}}
\newcommand{\setreg}[3]{\ensuremath{#3[#1 \mapsto #2]}}

% Memory Operations
\newcommand{\getat}[2]{\ensuremath{#2[#1]}}
\newcommand{\setat}[3]{\ensuremath{#3[#1 \mapsto \Allocated{#2}]}}
\newcommand{\allocat}[3]{\ensuremath{#3 \uplus \{#1 \mapsto \Allocated{#2}\}}}
\newcommand{\freeat}[2]{\ensuremath{#2[#1 \mapsto \Freed]}}

% State and Stack Constructors
\newcommand{\State}[3]{\textsf{State}\left(#1, #2, #3\right)}
\newcommand{\CallState}[2]{\textsf{CallState}\left(#1, #2\right)}
\newcommand{\ReturnState}[1]{\textsf{ReturnState}\left(#1\right)}
\newcommand{\Stackframe}[4]{\textsf{Stackframe}\left(#1, #2, #3, #4\right)}
\newcommand{\pcstate}{\textsf{pcstate}}
\newcommand{\rtlstate}{\textsf{rtl\_state}}

% RTL Instruction Set (single-line representations)
\newcommand{\Inop}[1]{\mathtt{nop} \to #1}
\newcommand{\Iop}[4]{#1 \coloneqq #2\; #3 \to #4}
\newcommand{\Iload}[3]{#1 \coloneqq {!#2} \to #3}
\newcommand{\Istore}[3]{{!#1} \coloneqq #2 \to #3}
\newcommand{\Ialloc}[2]{#1 \coloneqq \mathtt{alloc}() \to #2}
\newcommand{\Ifree}[2]{\mathtt{free}\; #1 \to #2}
\newcommand{\Icond}[3]{\mathtt{if}\; #1\; \mathtt{then}\; #2\; \mathtt{else}\; #3}
\newcommand{\Icall}[4]{#1 \coloneqq \mathtt{call}\; #2(#3) \to #4}
\newcommand{\Iret}[1]{\mathtt{ret}\; #1}
\newcommand{\Ireturn}[1]{\mathtt{ret}\; #1}

% RTL Helpers
\newcommand{\evalop}[2]{\textsf{eval\_op}\left(#1, #2\right)}
\newcommand{\findfun}[2]{\textsf{find\_fun}\left(#1, #2\right)}
\newcommand{\initregs}[2]{\textsf{init\_regs}\left(#1, #2\right)}

% Custom Operators
\makeatletter
\newcommand{\slantsqsubseteq}{\mathrel{\mathpalette\draw@slantsqsubseteq\relax}}
\newcommand{\draw@slantsqsubseteq}[2]{%
  \vcenter{\hbox{%
    \begin{tikzpicture}[
      baseline=0,
      line width=0.04em,
      line cap=round,
      line join=round,
      scale=0.08em
    ]
      % top arm down to base
      \draw (0.7em, 0.8em) -- (0em, 0.33em) -- (0.7em, 0.33em);
      % Underbar
      \draw (0em, 0.1em) -- (0.7em, 0.1em);
    \end{tikzpicture}%
  }}%
}
\makeatother


\newcommand{\Prop}{\textsf{Prop}} % Maybe \mathbb{P} ?
\newcommand{\reg}{\textsf{reg}}
\newcommand{\mem}{\textsf{memory}}
\newcommand{\regbank}{\textsf{regbank}}
\newcommand{\listT}[1]{\textsf{list}~{#1}}
\newcommand{\node}{\textsf{node}}
\newcommand{\instr}{\textsf{instr}}
\newcommand{\ident}{\textsf{ident}}
\newcommand{\pcatis}[3]{#1[#2] = \langle#3\rangle}
\newcommand{\true}{\textsf{true}}
\newcommand{\false}{\textsf{false}}

% Relational Separation Logic Macros
\newcommand{\rProp}{\textsf{rProp}}
\newcommand{\purel}[1]{\ulcorner #1 \urcorner}
\newcommand{\emp}{\textsf{emp}}
\providecommand{\sep}{\mathbin{\ast}}
\renewcommand{\sep}{\mathbin{\ast}}
\newcommand{\wand}{\mathrel{-\mkern-6mu\ast}}
\newcommand{\ent}{\vdash}
\newcommand{\boxm}{\Box}
\newcommand{\ptrt}[2]{#1 \to_{\text{t}} #2}
\newcommand{\ptrs}[2]{#1 \to_{\text{s}} #2}
\newcommand{\freet}[1]{\textsf{free}_{\text{t}}\left(#1\right)}
\newcommand{\frees}[1]{\textsf{free}_{\text{s}}\left(#1\right)}
\newcommand{\meminj}[2]{\textsf{meminj}_{#2}\left(#1\right)}
\newcommand{\sameargs}[3]{\textsf{same\_args}\left(#1, #2, #3\right)}
\newcommand{\samer}[3]{\textsf{samer}\left(#1, #2, #3\right)}
\newcommand{\hoareq}[4]{\{#1\}\; #2 \lesssim #3\; \{#4\}}
\newcommand{\bigast}{\mathop{\raisebox{-2pt}{\Large$\ast$}}\displaylimits}
